\documentclass[12pt]{article}
\usepackage[utf8]{inputenc}
\usepackage[T1]{fontenc}
\usepackage{hyperref}
\title{Yout Wam Language Resources for K-12 Curriculum Development}
\author{Compiled by Claude Sonnet 5.5 (claude-sonnet-5-5) for Andr\'as Kornai}
\date{}
\begin{document}
\maketitle

\section{Introduction}
Yout Wam (ISO 639-3: ytw; Glottocode yout1234) is a small Trans-New Guinea language of the Finisterre Mountains, Madang Province, Papua New Guinea. Search results place it in the Finisterre-Huon branch, Yupna group, close to Nankina Wam, Domung Meh and Yupno Gen. Population figures are small and disagree: Global Recordings Network gives about 270 and Joshua Project about 500 for the people group. Very little educational material was found.

\section{Primary Educational Resources}
No primers, textbooks or school programmes for Yout Wam were found in search results.

\section{Linguistic and Reference Materials}
The only reference-level pages found are catalogue entries.
\begin{itemize}
\item Crossgram (CLLD) language page for Yout Wam (yout1234): \url{https://crossgram.clld.org/languages/yout1234}
\item Wikipedia ISO 639 redirect for ytw, which leads to the article on the related Domung language. Domung's speaker figures do not apply to Yout Wam: \url{https://wikipedia.com/wiki/ISO_639:ytw}
\end{itemize}

\section{Religious and Community Resources}
\begin{itemize}
\item Joshua Project people-group profile for the Yout Wam. It lists Scripture status as ``Translation Started'': \url{https://joshuaproject.net/index.php/people_groups/20591}
\item Joshua Project language page for ytw: \url{https://joshuaproject.net/languages/ytw}
\end{itemize}

\section{Digital and Audio Resources}
Global Recordings Network has a language page for ytw, but it reports that no recordings are currently available: \url{https://globalrecordings.net/en/language/ytw}

\section{Curriculum Development Context}
Search results show no documented school use of Yout Wam. A curriculum would first need community-led documentation of the language, an orthography and a basic word list. Fieldwork and partnership with the community and local education authorities in Madang Province would be required. Nothing found here is a ready-made classroom resource.
\end{document}